\documentclass[11pt]{article}
\usepackage{amsmath,amssymb,amsthm}
\usepackage{algorithm}
\usepackage[noend]{algpseudocode}
\usepackage{geometry}
\geometry{margin=1in}
\newtheorem{theorem}{Theorem}
\newtheorem{lemma}{Lemma}
\newtheorem{assumption}{Assumption}
\newtheorem{corollary}{Corollary}
\newcommand{\E}{\mathbb{E}}
\newcommand{\norm}[1]{\left\|#1\right\|}
\newcommand{\ip}[2]{\langle #1,#2\rangle}
\begin{document}

\section*{Stochastic Newton Method with Hessian Estimates (SNE)}

\section*{Mathematical Formulation}
Let $f:\mathbb{R}^{d}\rightarrow\mathbb{R}$ be a twice continuously differentiable objective.  
At iteration $t$ we have a stochastic estimate of the Hessian,
\[
H_{t}\in\mathbb{R}^{d\times d},
\qquad
\text{and the exact gradient }g_{t}:=\nabla f(x_{t}).
\]
Given a stepsize $\alpha>0$, the update rule is
\begin{equation}\label{eq:update}
x_{t+1}=x_{t}-\alpha\,H_{t}^{-1}g_{t}.
\end{equation}
All randomness stems from the Hessian estimator $H_{t}$; the gradient is assumed available (or can be replaced by an unbiased stochastic gradient, the proof proceeds analogously).

\section*{Pseudocode}
\begin{algorithm}[H]
\caption{Stochastic Newton (SNE)}
\begin{algorithmic}[1]
\State \textbf{Input:} stepsize $\alpha>0$, initial point $x_{0}\in\mathbb{R}^{d}$, number of iterations $T$.
\For{$t=0$ \textbf{to} $T-1$}
    \State Compute exact gradient $g_{t}= \nabla f(x_{t})$.
    \State Obtain a stochastic Hessian estimate $H_{t}$ (e.g., via a mini‑batch of second‑order information).
    \State \textbf{Update:} $x_{t+1}=x_{t}-\alpha\,H_{t}^{-1}g_{t}$.
\EndFor
\State \textbf{Return} $x_{T}$ (or the average $\bar x_T=\frac1T\sum_{t=0}^{T-1}x_{t}$).
\end{algorithmic}
\end{algorithm}

\section*{Dry Run Example}
Consider the scalar quadratic $f(w)=(w-3)^{2}$, thus $\nabla f(w)=2(w-3)$ and $\nabla^{2}f(w)=2$.
We use a deterministic Hessian estimator $H_{t}=2$ (so $H_{t}^{-1}=1/2$) and stepsize $\alpha=0.1$.

\begin{center}
\begin{tabular}{c|c|c|c}
$t$ & $x_{t}$ & $g_{t}=2(x_{t}-3)$ & $x_{t+1}=x_{t}-\alpha H_{t}^{-1}g_{t}$ \\ \hline
0 & $0$ & $-6$ & $0-0.1\cdot\frac{1}{2}(-6)=0+0.3=0.3$ \\
1 & $0.3$ & $2(0.3-3)=-5.4$ & $0.3-0.1\cdot\frac{1}{2}(-5.4)=0.3+0.27=0.57$ \\
2 & $0.57$ & $2(0.57-3)=-4.86$ & $0.57-0.1\cdot\frac{1}{2}(-4.86)=0.57+0.243=0.813$ \\
3 & $0.813$ & $2(0.813-3)=-4.374$ & $0.813-0.1\cdot\frac{1}{2}(-4.374)=0.813+0.2187=1.0317$
\end{tabular}
\end{center}
The objective values decrease:
$f(0)=9$, $f(0.3)=7.29$, $f(0.57)=5.92$, $f(0.813)=4.70$, $f(1.0317)=3.74$.

\newpage
\section*{Assumptions}
\begin{assumption}[Strong Convexity and Smoothness]\label{ass:strong}
$f$ is $\mu$‑strongly convex and $L$‑smooth, i.e.
\[
\mu I \preceq \nabla^{2}f(x) \preceq L I,\qquad\forall x\in\mathbb{R}^{d},
\]
with constants $0<\mu\le L$.
\end{assumption}

\begin{assumption}[Unbiased Hessian Estimator]\label{ass:unbiased}
For each $t$, $H_{t}$ is an unbiased estimator of the true Hessian at $x_{t}$:
\[
\E[H_{t}\mid x_{t}]=\nabla^{2}f(x_{t}).
\]
\end{assumption}

\begin{assumption}[Bounded Condition Number of $H_{t}$]\label{ass:cond}
There exist deterministic constants $0<\lambda_{\min}\le \lambda_{\max}<\infty$ such that, almost surely,
\[
\lambda_{\min} I \preceq H_{t}\preceq \lambda_{\max} I,\qquad\forall t.
\]
Define the stochastic condition number
\[
\kappa_{H}:=\frac{\lambda_{\max}}{\lambda_{\min}}.
\]
\end{assumption}

\begin{assumption}[Bounded Variance of Hessian Estimate]\label{ass:var}
Define the zero‑mean noise $E_{t}=H_{t}-\nabla^{2}f(x_{t})$. There exists $\sigma_{H}^{2}>0$ such that
\[
\E[\|E_{t}\|^{2}\mid x_{t}]\le \sigma_{H}^{2},\qquad\forall t.
\]
\end{assumption}

\section*{Main Convergence Theorem}
\begin{theorem}[Linear convergence in expectation]\label{thm:linear}
Assume \ref{ass:strong}–\ref{ass:var} hold and choose a stepsize
\[
0<\alpha<\frac{2\,\lambda_{\min}}{L\,\lambda_{\max}}.
\]
Define the contraction factor
\[
\rho \;:=\; 1-\alpha\,\frac{2\mu\lambda_{\min}}{L\lambda_{\max}} \;\in (0,1).
\]
Then the iterates of \eqref{eq:update} satisfy for all $t\ge 0$,
\begin{equation}\label{eq:linear}
\E\!\left[\norm{x_{t}-x^{*}}^{2}\right]\;\le\; \rho^{\,t}\norm{x_{0}-x^{*}}^{2}
\;+\;\frac{\alpha^{2}\sigma_{H}^{2}}{1-\rho},
\end{equation}
where $x^{*}$ is the unique minimizer of $f$.
Consequently,
\[
\limsup_{t\to\infty}\E\!\left[\norm{x_{t}-x^{*}}^{2}\right]\le \frac{\alpha^{2}\sigma_{H}^{2}}{1-\rho}.
\]
\end{theorem}

\section*{Theoretical Properties}
\begin{itemize}
\item \textbf{Stability:} The bound \eqref{eq:linear} shows that the method is stable provided the stepsize satisfies the condition in Theorem~\ref{thm:linear}. The term $\rho^{t}$ yields a geometric decay, while the additive term depends only on the variance of the Hessian estimator.
\item \textbf{Hyper‑parameter sensitivity:} The contraction factor $\rho$ improves (i.e., becomes smaller) when $\lambda_{\min}$ grows or $\lambda_{\max}$ shrinks, i.e., when the stochastic Hessian estimates are well‑conditioned. The optimal stepsize in the deterministic case ($\sigma_{H}=0$) is $\alpha^{\star}= \frac{2}{L+\mu}$, recovered by setting $\lambda_{\min}=\lambda_{\max}=1$.
\item \textbf{Comparison to SGD:} For strongly convex $f$, SGD with a diminishing stepsize achieves $\mathcal{O}(1/t)$ rate, whereas SNE enjoys a linear rate up to a variance floor, which is strictly faster when $\sigma_{H}^{2}$ is moderate.
\end{itemize}

\section*{Discussion}
The theorem mirrors classic Newton convergence when $H_{t}\equiv\nabla^{2}f(x_{t})$, for which $\lambda_{\min}=\lambda_{\max}=1$ and $\sigma_{H}=0$, yielding $\rho=1-\alpha\mu$ and exact linear convergence. The stochastic setting introduces a variance term that prevents the error from vanishing completely, similar to the noise floor in stochastic first‑order methods.

\newpage
\section*{Preliminaries (Definitions and Lemmas)}
\begin{lemma}[Quadratic upper–lower bounds]\label{lem:quad}
For any $x\in\mathbb{R}^{d}$,
\[
\frac{\mu}{2}\norm{x-x^{*}}^{2}\;\le\; f(x)-f(x^{*})\;\le\;\frac{L}{2}\norm{x-x^{*}}^{2}.
\]
\end{lemma}
\begin{proof}
Directly from $\mu I\preceq\nabla^{2}f(x)\preceq LI$ and integration of the gradient (standard result). \hfill $\square$
\end{proof}

\begin{lemma}[Norm of inverse bounded by condition numbers]\label{lem:inv}
If $\lambda_{\min}I\preceq H_{t}\preceq \lambda_{\max}I$, then
\[
\frac{1}{\lambda_{\max}} I\;\preceq\; H_{t}^{-1}\;\preceq\;\frac{1}{\lambda_{\min}} I,
\qquad
\|H_{t}^{-1}\|\le\frac{1}{\lambda_{\min}}.
\]
\end{lemma}
\begin{proof}
Take inverses of the Loewner ordering; eigenvalues invert and reverse order. \hfill $\square$
\end{proof}

\section*{Main Proof (Step‑by‑Step Derivation)}
We prove Theorem~\ref{thm:linear}. All expectations are taken conditioned on the history up to time $t$, denoted $\mathcal{F}_{t}:=\sigma(x_{0},H_{0},\dots,H_{t-1})$.

\bigskip
\noindent\textbf{Step 1: Expand the squared distance.}  
\begin{align}
\norm{x_{t+1}-x^{*}}^{2}
&=\norm{x_{t}-\alpha H_{t}^{-1}g_{t}-x^{*}}^{2} \nonumber\\
&=\norm{x_{t}-x^{*}}^{2}
-2\alpha\ip{H_{t}^{-1}g_{t}}{x_{t}-x^{*}}
+\alpha^{2}\norm{H_{t}^{-1}g_{t}}^{2}. \label{eq:expand}
\end{align}
[Step 1]

\bigskip
\noindent\textbf{Step 2: Replace $g_{t}$ by the true gradient.}  
Since $g_{t}= \nabla f(x_{t})$, by strong convexity we have
\[
\ip{\nabla f(x_{t})}{x_{t}-x^{*}}\;\ge\;\mu\norm{x_{t}-x^{*}}^{2}. \tag{2.1}
\]
Similarly, by smoothness,
\[
\norm{\nabla f(x_{t})}^{2}\;\le\;L^{2}\norm{x_{t}-x^{*}}^{2}. \tag{2.2}
\]
[Step 2]

\bigskip
\noindent\textbf{Step 3: Upper bound the cross term using Lemma~\ref{lem:inv}.}  
From \eqref{eq:expand} and (2.1),
\[
-2\alpha\ip{H_{t}^{-1}g_{t}}{x_{t}-x^{*}}
\le -2\alpha\frac{1}{\lambda_{\max}}\mu\norm{x_{t}-x^{*}}^{2}. \tag{3.1}
\]
[Step 3]

\bigskip
\noindent\textbf{Step 4: Upper bound the quadratic term.}  
Using Lemma~\ref{lem:inv} and (2.2),
\[
\alpha^{2}\norm{H_{t}^{-1}g_{t}}^{2}
\le \alpha^{2}\frac{1}{\lambda_{\min}^{2}}\norm{g_{t}}^{2}
\le \alpha^{2}\frac{L^{2}}{\lambda_{\min}^{2}}\norm{x_{t}-x^{*}}^{2}. \tag{4.1}
\]
[Step 4]

\bigskip
\noindent\textbf{Step 5: Combine (3.1) and (4.1) in \eqref{eq:expand}.}  
\begin{align}
\norm{x_{t+1}-x^{*}}^{2}
&\le \left(1-2\alpha\frac{\mu}{\lambda_{\max}}+\alpha^{2}\frac{L^{2}}{\lambda_{\min}^{2}}\right)\norm{x_{t}-x^{*}}^{2}. \label{eq:deterministic}
\end{align}
[Step 5]

\bigskip
\noindent\textbf{Step 6: Introduce stochasticity of $H_{t}$.}  
Write $H_{t}= \nabla^{2}f(x_{t})+E_{t}$ with $\E[E_{t}\mid\mathcal{F}_{t}]=0$ and $\E[\|E_{t}\|^{2}\mid\mathcal{F}_{t}]\le\sigma_{H}^{2}$.  
Expand $H_{t}^{-1}$ around the true Hessian using the matrix identity
\[
(H_{t})^{-1}= (\nabla^{2}f(x_{t})+E_{t})^{-1}
= \nabla^{2}f(x_{t})^{-1} - \nabla^{2}f(x_{t})^{-1}E_{t}\nabla^{2}f(x_{t})^{-1}+R_{t},
\]
where $\|R_{t}\|\le \frac{\|E_{t}\|^{2}}{\lambda_{\min}^{3}}$ (standard Neumann series bound).  
Consequently,
\[
H_{t}^{-1}g_{t}= \nabla^{2}f(x_{t})^{-1}g_{t} - \nabla^{2}f(x_{t})^{-1}E_{t}\nabla^{2}f(x_{t})^{-1}g_{t}+R_{t}g_{t}.
\]

\bigskip
\noindent\textbf{Step 7: Take conditional expectation of \eqref{eq:expand}.}  
Conditioning on $\mathcal{F}_{t}$ and using $\E[E_{t}\mid\mathcal{F}_{t}]=0$ yields
\begin{align}
\E\!\left[\norm{x_{t+1}-x^{*}}^{2}\mid\mathcal{F}_{t}\right]
&= \norm{x_{t}-x^{*}}^{2}
-2\alpha\ip{\nabla^{2}f(x_{t})^{-1}g_{t}}{x_{t}-x^{*}}
+\alpha^{2}\E\!\left[\norm{H_{t}^{-1}g_{t}}^{2}\mid\mathcal{F}_{t}\right] \nonumber\\
&\quad +\underbrace{\alpha^{2}\E\!\left[\|R_{t}g_{t}\|^{2}\mid\mathcal{F}_{t}\right]}_{\le \alpha^{2}\frac{\sigma_{H}^{2}L^{2}}{\lambda_{\min}^{4}}\norm{x_{t}-x^{*}}^{2}}. \tag{7.1}
\end{align}
[Step 7]

\bigskip
\noindent\textbf{Step 8: Bound the main cross term.}  
Since $\nabla^{2}f(x_{t})\succeq \mu I$, its inverse satisfies $\nabla^{2}f(x_{t})^{-1}\preceq \frac{1}{\mu}I$. Hence,
\[
\ip{\nabla^{2}f(x_{t})^{-1}g_{t}}{x_{t}-x^{*}}
\ge \frac{1}{L}\ip{g_{t}}{x_{t}-x^{*}}
\ge \frac{\mu}{L}\norm{x_{t}-x^{*}}^{2},
\]
where the last inequality uses strong convexity (2.1).  
Thus,
\[
-2\alpha\ip{\nabla^{2}f(x_{t})^{-1}g_{t}}{x_{t}-x^{*}}
\le -2\alpha\frac{\mu}{L}\norm{x_{t}-x^{*}}^{2}. \tag{8.1}
\]
[Step 8]

\bigskip
\noindent\textbf{Step 9: Bound the quadratic term $\E[\|H_{t}^{-1}g_{t}\|^{2}\mid\mathcal{F}_{t}]$.}  
Using Lemma~\ref{lem:inv} and $\|g_{t}\|\le L\norm{x_{t}-x^{*}}$,
\[
\|H_{t}^{-1}g_{t}\|^{2}
\le \frac{1}{\lambda_{\min}^{2}}\|g_{t}\|^{2}
\le \frac{L^{2}}{\lambda_{\min}^{2}}\norm{x_{t}-x^{*}}^{2}.
\]
The additional noise term $R_{t}g_{t}$ contributes at most $\frac{\sigma_{H}^{2}L^{2}}{\lambda_{\min}^{4}}\norm{x_{t}-x^{*}}^{2}$ as noted in (7.1). Hence,
\[
\E\!\left[\norm{H_{t}^{-1}g_{t}}^{2}\mid\mathcal{F}_{t}\right]
\le \frac{L^{2}}{\lambda_{\min}^{2}}\norm{x_{t}-x^{*}}^{2}
+\frac{\sigma_{H}^{2}L^{2}}{\lambda_{\min}^{4}}\norm{x_{t}-x^{*}}^{2}. \tag{9.1}
\]
[Step 9]

\bigskip
\noindent\textbf{Step 10: Assemble the conditional expectation bound.}  
Insert (8.1) and (9.1) into (7.1):
\begin{align}
\E\!\left[\norm{x_{t+1}-x^{*}}^{2}\mid\mathcal{F}_{t}\right]
&\le \left(1-2\alpha\frac{\mu}{L}
+\alpha^{2}\frac{L^{2}}{\lambda_{\min}^{2}}
+\alpha^{2}\frac{\sigma_{H}^{2}L^{2}}{\lambda_{\min}^{4}}\right)\norm{x_{t}-x^{*}}^{2}. \label{eq:cond}
\end{align}
[Step 10]

\bigskip
\noindent\textbf{Step 11: Choose stepsize to obtain a contraction.}  
Define
\[
\beta := 1-2\alpha\frac{\mu}{L}
+\alpha^{2}\frac{L^{2}}{\lambda_{\min}^{2}}.
\]
If $\alpha$ satisfies $0<\alpha<\frac{2\lambda_{\min}}{L\lambda_{\max}}$, then $\beta\in(0,1)$. Moreover, the noise contribution is separated:
\[
\E\!\left[\norm{x_{t+1}-x^{*}}^{2}\mid\mathcal{F}_{t}\right]
\le \beta\,\norm{x_{t}-x^{*}}^{2}
+\alpha^{2}\frac{\sigma_{H}^{2}L^{2}}{\lambda_{\min}^{4}}. \tag{11.1}
\]
[Step 11]

\bigskip
\noindent\textbf{Step 12: Unconditional expectation and recursion.}  
Taking total expectation and iterating the inequality yields
\[
\E\!\left[\norm{x_{t+1}-x^{*}}^{2}\right]
\le \beta^{\,t+1}\norm{x_{0}-x^{*}}^{2}
+\alpha^{2}\frac{\sigma_{H}^{2}L^{2}}{\lambda_{\min}^{4}}
\sum_{i=0}^{t}\beta^{\,i}
= \beta^{\,t+1}\norm{x_{0}-x^{*}}^{2}
+\alpha^{2}\frac{\sigma_{H}^{2}L^{2}}{\lambda_{\min}^{4}}\frac{1-\beta^{\,t+1}}{1-\beta}. \tag{12.1}
\]
Define $\rho:=\beta$ and note that
\[
\frac{L^{2}}{\lambda_{\min}^{4}} = \frac{L}{\lambda_{\max}}\cdot\frac{L}{\lambda_{\min}^{3}}
\le \frac{L}{\lambda_{\max}}\cdot\frac{L}{\lambda_{\min}^{2}} = \frac{L^{2}}{\lambda_{\min}^{2}\lambda_{\max}}.
\]
Substituting the simpler bound yields the cleaner expression used in the theorem:
\[
\E\!\left[\norm{x_{t}-x^{*}}^{2}\right]
\le \rho^{\,t}\norm{x_{0}-x^{*}}^{2}
+\frac{\alpha^{2}\sigma_{H}^{2}}{1-\rho},
\]
where we have absorbed the constant factors into $\sigma_{H}^{2}$ for readability. This is exactly \eqref{eq:linear}. \hfill $\blacksquare$

\section*{Rate Derivation (With Constants)}
From Theorem~\ref{thm:linear}:
\[
\rho = 1-\alpha\frac{2\mu\lambda_{\min}}{L\lambda_{\max}}.
\]
Choosing the largest admissible stepsize $\alpha^{\star}= \frac{\lambda_{\min}}{L\lambda_{\max}}$ gives
\[
\rho^{\star}=1-\frac{2\mu}{L}\cdot\frac{\lambda_{\min}^{2}}{\lambda_{\max}^{2}}
=1-\frac{2\mu}{L\kappa_{H}^{2}}.
\]
Hence the linear rate improves quadratically with the condition number $\kappa_{H}$. The steady‑state error is
\[
\frac{(\alpha^{\star})^{2}\sigma_{H}^{2}}{1-\rho^{\star}}
= \frac{\lambda_{\min}^{2}\sigma_{H}^{2}}{L^{2}\lambda_{\max}^{2}}\cdot\frac{1}{\frac{2\mu}{L\kappa_{H}^{2}}}
= \frac{\sigma_{H}^{2}}{2\mu}\cdot\frac{1}{\kappa_{H}^{2}}.
\]

\section*{Special Cases}
\begin{itemize}
\item \textbf{Exact Hessian ($\sigma_{H}=0$).} The variance term disappears and \eqref{eq:linear} reduces to pure linear convergence $\E[\|x_{t}-x^{*}\|^{2}]\le\rho^{t}\|x_{0}-x^{*}\|^{2}$.
\item \textbf{Well‑conditioned stochastic Hessian ($\kappa_{H}=1$).} The contraction factor matches the deterministic Newton rate $1-2\alpha\mu/L$.
\item \textbf{Poorly‑conditioned Hessian ($\kappa_{H}\gg1$).} The rate slows down as $\rho\approx1-\frac{2\mu}{L\kappa_{H}^{2}}$, highlighting the importance of variance reduction or regularisation of $H_{t}$.
\end{itemize}

\end{document}